% ---
% title: Recent Image Compression：演讲提纲
% date: 2026-10-14
% tag: Slides
% excerpt: 配合幻灯片 Recent Image Compression 的脱稿提纲。先记住一句主线和四字口诀“省、真、轻、机”，再按页记每页要讲的点和接下一页的过渡句；论文页的图是原论文的 pipeline，指着图从左往右讲；数字都在幻灯片上，不用背。附时间检查点、卡住时的办法和可能的提问。
% ---
\definecolor{sig}{HTML}{1F5FAD}
\definecolor{hot}{HTML}{C0392B}
\definecolor{accent}{HTML}{C77C02}
\definecolor{soft}{HTML}{8C8C8C}
\definecolor{video}{HTML}{2E7D6B}

\begin{abstract}
这是幻灯片 \href{https://mejistus.github.io/\#post=Slides/Recent-Image-Compression.tex}{Recent Image Compression} 的提纲，给没有提词器、要脱稿讲完全场的情况用。
要背的只有三样：\textbf{第~\ref{sec:spine}~节}的一句主线和四字口诀；\textbf{第~\ref{sec:pages}~节}每页的“接”（过渡句）；以及每篇论文都按 Problem、Method、Result 三拍讲。
论文页上方是原论文的 pipeline 图，讲 Method 时指着图从左往右讲；数字、会议、年份都印在幻灯片上，指着读，不用背。
准备顺序：先能不看稿说出主线和四支的表；再对着幻灯片把第~\ref{sec:pages}~节出声讲两遍；最后只看幻灯片讲一遍，计时。
\end{abstract}

%======================================================================
\section{主线：一句话，四个字}\label{sec:spine}

\textbf{一句话：}学习型编码在 2020 年追平了 VVC；之后分成四支，每支改同一条流水线上的一个地方，也都把代价挪到了别处；2025 年的 JPEG AI 把它们收进了一个标准。今天每支只讲 2025 年以后发表的两篇。

\textbf{口诀：省、真、轻、机。}四支的顺序、颜色和幻灯片第 2 页的四个圆圈一致：

\begin{center}
\begin{tabular}{lllll}
\toprule
口诀 & 分支 & 问的问题 & 两篇代表作 & 代价挪到哪 \\
\midrule
\textcolor{sig}{\textbf{省}} & 1 Rate & 同样误差，能少几个比特？ & DCAE, HPCM & 解码计算更多 \\
\textcolor{accent}{\textbf{真}} & 2 Perception & 比特极少时，图还像真的吗？ & StableCodec, OneDC & 细节是编出来的 \\
\textcolor{video}{\textbf{轻}} & 3 Complexity & 解码能便宜到 CPU 吗？ & LotteryCodec, Good-Cheap-Fast & 编码变慢 \\
\textcolor{hot}{\textbf{机}} & 4 Receiver & 看图的是大模型，怎么压？ & Transform-neck, CoTAM & 绑定一个视觉编码器 \\
\bottomrule
\end{tabular}
\end{center}

\textbf{每支两篇怎么记：}第一支，DCAE 往图外看，HPCM 在图内看得更远；第二支，StableCodec 让一步去噪就够，OneDC 换掉文字提示；第三支，LotteryCodec 不传权重，Good-Cheap-Fast 让小解码器也能“真”；第四支，transform-neck 不要像素，CoTAM 改比特花在哪。

\textbf{每支同一个节奏：}先一页“这一支在问什么”，再两页论文。论文页上的三行粗体就是提示：\textbf{Problem}（它要解决什么）$\to$ \textbf{Method}（怎么做，指着图讲）$\to$ \textbf{Result}（念数字）。最后补半句“代价是……”。

%======================================================================
\section{逐页提纲}\label{sec:pages}

每页三行：\textbf{说}是要讲到的点，按顺序；\textbf{接}是翻页前的最后一句，它引出下一页的标题。括号里是这一页的时长和讲完时的累计时间。

\subsection*{开场}

\textbf{第 1 页\quad 标题}（0:15，0:15）
\begin{itemize}
  \item 说：这是上次 From Samples to Bits 的续篇。上次讲到学习型编码追平了 VVC；今天讲之后它往哪几个方向走，每个方向讲两篇 2025 年以后的工作。
  \item 接：“先回顾上次停在哪。”
\end{itemize}

\textbf{第 2 页\quad Where we left off}（1:00，1:15）
\begin{itemize}
  \item 说：图从左往右读：$x$ 经过变换 $g_a$、量化 $Q$、熵编码 AE 变成比特；解码端反过来，AD、$g_s$，得到 $\hat x$。上面的熵模型 $p$ 告诉编解码器每个符号有多可能。
  \item 说：公式 $R+\lambda D$：$R$ 是比特数，$D$ 是误差，$\lambda$ 调两者的比重。2020 年这一套追平了 VVC。
  \item 说：之后分成四支，四个圆圈就是各自改的地方：1 改熵模型和变换；2 改解码器，让它会“画”；3 把解码器变小；4 换掉接收者。
  \item 接：“每一支问一个问题，下一页是总表。”
\end{itemize}

\textbf{第 3 页\quad Four branches, eight papers, one standard}（0:45，2:00）
\begin{itemize}
  \item 说：按行读，每支一个问题、两篇代表作、发表在哪；八篇都是 2025、2026 年正式发表的，更早的工作只在每支开头当背景。最后一行是 JPEG AI 标准。
  \item 说：底下的尺子 BD-rate：同样质量下比特变化的百分比，负数是省，$-10\%$ 就是省一成。
  \item 接：“先看第一支：同样的误差，能不能少几个比特。”
\end{itemize}

\subsection*{\textcolor{sig}{第一支：省}（2:00 开始）}

\textbf{第 4 页\quad Branch 1}（1:00，3:00）
\begin{itemize}
  \item 说：一个符号要花 $-\log_2 p$ 个比特，预测越准，比特越少。
  \item 说：能改的零件有两个：变换 $g_a$，要同时抓住局部纹理和长程结构；熵模型 $p$，要用解码端已有的信息去预测。
  \item 说（指图）：$\hat y$ 分成几片按顺序解码；橙色这一片用已经解出的片和超先验 $\hat z$ 来预测。
  \item 说：2025 年的两篇都改熵模型：一篇往图外看，一篇在图内看得更远。
  \item 接：“先看往图外看的：DCAE。”
\end{itemize}

\textbf{第 5 页\quad DCAE}（1:10，4:10）
\begin{itemize}
  \item Problem：熵模型只看这张图自己，可边缘、纹理在整个训练集里反复出现。
  \item Method（指图）：左上，$\hat y$ 分成五片按顺序解，每片一个模块 $e_i$；蓝框是放大的 $e_i$。右下橙框是关键：一本字典，128 条、每条 640 维，存在模型里，不占比特；当前特征当 Q，字典当 K 和 V，用交叉注意力去查。
  \item Result：比 VTM-12.1 省 17.0\%、21.1\%、19.7\%；Kodak 一张图 193 毫秒，MLIC++ 要 772 毫秒，快四倍。
  \item 接：“DCAE 往图外看；HPCM 在图内看得更远。”
\end{itemize}

\textbf{第 6 页\quad HPCM}（1:10，5:20）
\begin{itemize}
  \item Problem：上下文看得越远越有用，可看得远很贵；而且每一步只用自己的上下文，前几步算出来的东西就丢了。
  \item Method（指图，从左往右）：$\hat y$ 按位置抽成三个尺度，绿色最小（$\frac14$）、蓝色 $\frac12$、橙色全尺寸。先编最小的 2 步，再编 $\frac12$ 的 3 步，最后全尺寸 6 步，共 11 步，格子里的数字就是第几步。先编小尺度，等于先看远处。底下黄色箭头是渐进式上下文融合：每一步的上下文用交叉注意力传给下一步，不丢。
  \item Result：比 VTM-22.0 省 19.2\%、22.2\%、18.4\%；解码 90 毫秒，MLIC++ 226 毫秒，计算量差不多。
  \item 接：“第一支一直在省比特。可比特很少的时候，还有一个问题：图会糊。”
\end{itemize}

\subsection*{\textcolor{accent}{第二支：真}（5:20 开始）}

\textbf{第 7 页\quad Branch 2}（1:00，6:20）
\begin{itemize}
  \item 说：低码率时，MSE 最优的解码器输出的是条件平均，平均出来就是糊的。
  \item 说（指红线）：码率固定时，越真实，误差越大。线的一头是“MSE 最优、糊”，另一头是“完全真实”；Blau 和 Michaeli 证明，完全真实最多让 MSE 翻倍。
  \item 说：所以只传结构，纹理让扩散模型去生成；2025 年的变化是扩散只走一步。（指红字）代价：生成的不是传过来的。
  \item 接：“第一篇：StableCodec，怎么让一步就够。”
\end{itemize}

\textbf{第 8 页\quad StableCodec}（1:10，7:30）
\begin{itemize}
  \item Problem：扩散要 20 到 50 步，而且会偏离原图。
  \item Method（指图，从左往右）：编码端，$x$ 先过 Stable Diffusion 的编码器得到潜变量 $l$，再用中间橙框里压得很狠的潜变量编码器变成比特，空间上压 64 倍。解码端，$g_s$ 解出的不是干净的潜变量，而是带噪的 $l_T$，正好让蓝框里一步的 SD-Turbo 去噪，提示词是固定的一句“高分辨率、8K”。上面红线是辅助分支 $\mathcal D_{\mathrm{Aux}}$，从 $\hat y$ 直接带结构，加到 $l_0$ 上，防止偏离原图。
  \item Result（指表）：1 步、0.33 秒，和 GAN 编码器 MS-ILLM 一样快；PerCo 20 步 3.74 秒，DiffEIC 50 步 7.42 秒。0.005 bpp 时 FID、DISTS 最好。
  \item 接：“StableCodec 的提示词是固定的一句话。OneDC 问：一步去噪，该拿什么当提示？”
\end{itemize}

\textbf{第 9 页\quad OneDC}（1:10，8:40）
\begin{itemize}
  \item Problem：一步生成更依赖引导；文字说不清东西在哪、长什么样。
  \item Method（指图）：左边黄框是普通的超先验编码器，$\hat y$ 是主码流；下面的 $z$ 用 FSQ 量化成离散 token $\hat z$，只占 0.0034 bpp。关键在中间：$\hat z$ 经过 $h_{sem}$ 变成语义向量 $c$，代替文字，送进 U-Net 的交叉注意力。右边红框是 SD 1.5 的 U-Net 加 LoRA，一步生成。训练时还让 $\hat z$ 向 MaskGIT 的图像 tokenizer 学语义。
  \item Result：MS-COCO 上同样 FID，比 PerCo 少 39.6\% 比特；解码一张 $1024\times1024$ 的图 0.34 秒，PerCo 要 8.8 秒（A100），快 20 多倍。
  \item 接：“第二支的解码器有近十亿参数。第三支反过来问：解码器能不能小到 CPU 就能跑。”
\end{itemize}

\subsection*{\textcolor{video}{第三支：轻}（8:40 开始）}

\textbf{第 10 页\quad Branch 3}（1:00，9:40）
\begin{itemize}
  \item Problem：一个解码器要应付所有图，所以很大，JPEG AI 的每像素要 1.3 万到 21.5 万次乘加。
  \item 说：想法是过拟合，只为这一张图训练一个小模型，把模型和图一起发出去。
  \item 说（指图）：这是 Cool-Chic 和 C3 的解码过程。A，熵模型是三层 $1\times1$ 卷积的小网络，看已解码的邻居，给出 $\mu$ 和 $\sigma$；B，几层不同分辨率的潜变量上采样后叠在一起，过一个小的合成网络出 RGB。底下的比特流里，除了潜变量，还有这两个小网络的权重。
  \item 说：这类方法每像素约 3k 次乘加；它们是背景，今天的两篇在它们之上。（指红字）代价挪到了编码端：每张图都是一次训练。
  \item 接：“比特流里那段权重，码率低时很占地方。LotteryCodec 干脆不传权重。”
\end{itemize}

\textbf{第 11 页\quad LotteryCodec}（1:10，10:50）
\begin{itemize}
  \item Problem：过拟合编码要把解码器的权重一起发；码率低时，权重占了大头。
  \item Method（指图）：上半是编码：在一个随机初始化的大网络里找一张二值掩码，就是剪刀剪掉的那些边，再加一组调制向量，熵编码后发出去。下半是解码：用同一个随机种子把同一个随机网络重新生成出来，套上掩码，就是这张图的解码器。叫“彩票”，是因为随机网络里本来就藏着一个好用的子网络，只要把它找出来。
  \item Result：Kodak 和 CLIC 上都超过 VTM-19.1，CLIC 上最多省 9.8\%；每像素约 4k 次乘加；掩码比例可调，用乘加换比特。
  \item 接：“LotteryCodec 仍按 MSE 训练，码率低时图会发平。下一篇让小解码器也能‘真’。”
\end{itemize}

\textbf{第 12 页\quad Good, Cheap, and Fast}（1:10，12:00）
\begin{itemize}
  \item Problem：HiFiC、扩散这些生成式编码器，解码每像素要 10 万到 100 万次乘加；便宜的过拟合编码器又按 MSE 训练，纹理发平。
  \item Method（指图）：还是 C3，只换了损失，叫 Wasserstein Distortion。A，用 VGG 提特征；B，在每个位置的邻域里算特征的均值、方差这些统计量；C，比两张图的局部统计量，而不是逐像素比。意思是纹理只要统计上像就行，和人眼余光的感知一致。另外，解码器还拿到一份由固定种子生成的噪声，不占比特，用来生成纹理。
  \item Result：人眼评测，CLIC 2020、0.15 bpp，和 HiFiC 差不多（Elo 2187 对 2224）；解码每像素 3.1k 次乘加，HiFiC 要 60.9 万，不到 1\%。
  \item 接：“这两篇都用了一个共享的随机种子，不占比特。前三支的接收者都是人眼。第四支问：如果看图的是多模态大模型呢？”
\end{itemize}

\subsection*{\textcolor{hot}{第四支：机}（12:00 开始）}

\textbf{第 13 页\quad Branch 4}（0:50，12:50）
\begin{itemize}
  \item 说（指图左边三格）：(a) 用为人眼训练的编解码器，解出图再给 MLLM，任务性能不好；(b)(c) 把任务损失从 MLLM 一路反传回编码器，性能好，但 MLLM 太大，训不动。右边 (d) 就是下一页的 transform-neck。
  \item 说：2025、2026 年这两篇的共同思路：只对齐 MLLM 的视觉编码器，也就是 CLIP，不碰 LLM。
  \item 接：“第一篇：transform-neck，连像素都不要了。”
\end{itemize}

\textbf{第 14 页\quad Transform-neck}（1:10，14:00）
\begin{itemize}
  \item Problem：低码率解出的图会拉低 MLLM 的准确率；LLM 太大，没法连着一起训练。
  \item Method（指图，从左往右）：终端只跑编码器；中间传比特；云端不解码成像素，transform-neck 直接把潜变量 $\hat y$ 变成 CLIP 的中间层输入，跳过 CLIP 的前两层。它只有 13M 参数，训练只用 CLIP 的特征，所以用同一个 CLIP 的 MLLM 都能直接用。
  \item Result：同样准确率省 60\% 到 80\% 的比特；云端从比特到 CLIP 的计算从每像素 1018k 次乘加降到 53k，少约 95\%。
  \item 接：“transform-neck 不要像素了。CoTAM 保留像素，改的是比特花在哪。”
\end{itemize}

\textbf{第 15 页\quad CoTAM}（1:10，15:10）
\begin{itemize}
  \item Problem：作者先分析：压缩失真对 MLLM 各层特征的影响不均匀。低层特征退化得慢，可把细节和上下文连起来的那类特征最先坏。
  \item Method（指图）：右边是推理：编码端跑 CLIP 的前三层，把注意力平均成一张重要性图，量化后告诉编码器哪里多给比特。解码端在解出的图之外，加一个只有一个 Transformer 块的 adapter，把潜变量的特征送进 MLLM。左边是训练：编解码器和 MLLM 都冻结，只训 adapter，损失在 CLIP 的好几层上一起算。
  \item Result：LLaVA-1.5、六个任务，同样分数下比 ELIC 少 36\% 比特；换 DCAE 当底座，少 31\%，DCAE 就是第一支讲的那篇。时间只多不到 6\%。
  \item 接：“四支讲完了。它们在哪里汇合？JPEG AI。”
\end{itemize}

\subsection*{收尾（15:10 开始）}

\textbf{第 16 页\quad JPEG AI}（1:00，16:10）
\begin{itemize}
  \item 说：第一个学习型图像标准，ISO/IEC 6048-1，2025 年发布。
  \item 说（指左图）：这是它的解码器，就是第 2 页那条流水线：码流经过算术解码器和超先验解码器得到 $\hat y$，再过合成变换出图；亮度、色度各一路。
  \item 说（按颜色读右边）：\textcolor{sig}{率}，比 VVC intra 省 14\% 到 27\%（7 个指标的平均）；\textcolor{video}{复杂度}，一个码流、三种解码器，设备自己挑；比特精确，熵解码全用整数，每台设备算出的概率都一样；\textcolor{hot}{接收者}，机器任务在范围内，但第一版面向人眼。
  \item 接：“最后用一张表把四支收起来。”
\end{itemize}

\textbf{第 17 页\quad Four branches in one table}（1:00，17:10）
\begin{itemize}
  \item 说：按行读，问题、方法、做到了什么。
  \item 说（全场结论，放慢）：每一支都只是把代价挪了个地方：更多的解码计算、编出来的细节、更慢的编码、绑定一个视觉编码器。
  \item 接：“还有哪些没解决？”
\end{itemize}

\textbf{第 18 页\quad Open problems}（0:50，18:00）
\begin{itemize}
  \item 说：五条，各一句：生成的细节能不能信；过拟合编码怎么变快；所有设备上比特精确地解码；一个码流同时给人和 MLLM；趋势是编码器越来越多地建在预训练模型上，今天八篇里有四篇直接用了 Stable Diffusion 或 CLIP。
  \item 接：“谢谢，欢迎提问。”（两页参考文献直接翻过，停在 Thank you 页。）
\end{itemize}

%======================================================================
\section{时间检查点}

\begin{center}
\begin{tabular}{lll}
\toprule
到这一页 & 应该在 & 慢了怎么办 \\
\midrule
第 7 页（分支 2） & 5:20 & 第 8 页只讲“带噪潜变量加一步去噪”，然后指表 \\
第 10 页（分支 3） & 8:40 & 第 11 页只讲“传掩码，不传权重”和结果 \\
第 13 页（分支 4） & 12:00 & 第 13 页一句带过，直接翻到第 14 页 \\
第 16 页（JPEG AI） & 15:10 & 第 18 页只说前两条 \\
第 21 页（结束） & 18:10 & \\
\bottomrule
\end{tabular}
\end{center}

%======================================================================
\section{卡住时}

\begin{itemize}
  \item \textbf{先念标题。}每页标题都是这一页的结论，念出来再解释。
  \item \textbf{论文页按三行粗体讲。}Problem、Method、Result 印在幻灯片上，顺着念就不会乱。
  \item \textbf{讲图从左往右。}输入在左，输出在右；先说输入是什么、输出是什么，再说中间哪一块是这篇新加的。
  \item \textbf{数字不背。}指着那一行念出来，听众也在看。
  \item \textbf{忘了过渡句，就用通用句：}“这一篇解决了……，代价是……；下一篇要问的是……”。“代价是什么”在第~\ref{sec:spine}~节的表里。
  \item \textbf{整段忘了，回到口诀：}省、真、轻、机，说出现在讲到第几个字，再接着讲。
\end{itemize}

%======================================================================
\section{可能被问到的}

\textbf{DCAE 和 HPCM 的数字能直接比吗？}
不能直接比。DCAE 对比的是 VTM-12.1，HPCM 对比的是 VTM-22.0，锚点版本不同。两篇的表里都有 MLIC++，Kodak 上分别是 $-15.1\%$ 和 $-15.15\%$，所以大致可以参照：DCAE $-17.0\%$，HPCM-Large $-19.2\%$。严格比较要在同一张表里。

\textbf{生成式压缩的细节能信吗？}
不能保证。纹理是生成的，不是传过来的，文字和人脸可能被编出来。所以医疗、取证一类场景不宜用，这也是开放问题的第一条。

\textbf{过拟合编码这么慢，有什么用？}
适合编码一次、解码很多次的场景，比如网页和应用里的图片：编码在服务器上慢慢做，用户的手机解码很便宜。今天这两篇都没有解决编码慢，这仍是开放问题。

\textbf{LotteryCodec 的“每像素约 4k 次乘加”是实测吗？}
不是。这是只算没被掩掉的权重的理论值，要用支持稀疏计算的库才能接近；论文也给了不剪枝时的上限，并用结构化剪枝测了实际时间。

\textbf{为什么选这几篇？}
八篇都在 2025、2026 年正式发表（CVPR、ICCV、NeurIPS、ICML、ICLR）。每支两篇从不同角度解同一个问题，见第~\ref{sec:spine}~节“每支两篇怎么记”。更早的 TCM、PerCo、Cool-Chic、C3 只在分支开头和对比表里当背景。

\textbf{JPEG AI 用的是哪一支？}
主要是第一支和第三支：码率比 VVC intra 低，又按复杂度分了三档解码器。第四支在它的范围内，但第一版面向人眼。
